\documentclass[12pt, a4paper]{article}
\usepackage[utf8]{inputenc}
\usepackage[T1]{fontenc}
\usepackage{graphicx}
\usepackage{float}
\usepackage[margin=1in]{geometry}

\title{\textbf{Part 1: Conceptual Database Design}}
\author{
    \begin{tabular}{c}
    Yüksel Ege Boyacı -- 2023400315 \\
    Tuğçe Yücel -- 2022400105
    \end{tabular}
}
\date{}

\begin{document}

\maketitle
\newpage
\tableofcontents
\newpage

\section{Entity-Relationship (ER) Diagram}

The conceptual database design was created using the standard Entity-Relationship (ER) notation described in the lectures. The diagram below models all entity sets, relationship sets, key constraints, participation constraints, and class hierarchies (Is-A relationships) as specified in the project requirements.

\begin{figure}[H]
    \centering
    % Note: You may need to adjust the path depending on where your main .tex file is compiled from.
    % Assuming main file is in reports/latex or similar:
    \includegraphics[width=\textwidth]{er_diagram_final.png}
    \caption{Conceptual ER Diagram for TransferDB}
    \label{fig:er_diagram}
\end{figure}

\newpage
\section{Legend}
The graphical notation used in the ER diagram follows these conventions:
\begin{itemize}
    \item \textbf{Blue rectangle:} Entity
    \item \textbf{Yellow oval:} Attribute
    \item \textbf{Green diamond:} Relationship
    \item \textbf{Red triangle:} ISA relationship
    \item \textbf{Underlined:} Primary Key
    \item \textbf{Bold arrow toward relation:} Total participation + Key constraint
    \item \textbf{Thin arrow toward relation:} Partial participation + Key constraint
    \item \textbf{Bold line toward relation:} Total participation + No key constraint
    \item \textbf{Thin line toward relation:} Partial participation + No key constraint
\end{itemize}

\newpage
\section{Constraints Not Expressible in ER Design}

While the ER diagram effectively captures the core entities and overall structural relationships, certain complex business rules, cross-table temporal validations, and numeric derivations cannot be natively expressed in standard ER diagramming notation. These constraints will be addressed at the logical layer via SQL \texttt{CHECK} constraints or database triggers (or application-level logic). The un-expressible constraints are as follows:

\begin{enumerate}
    \item (\texttt{stadium\_name}, \texttt{city}) must be UNIQUE.
    \item (\texttt{competition\_name}, \texttt{season}) must be UNIQUE.
    \item \texttt{attendance} $\le$ \texttt{stadium capacity}.
    \item No two matches within 120 minutes at the same stadium.
    \item No club in two matches within 120 minutes.
    \item No referee in two matches within 120 minutes.
    \item Max 11 starters per club per match.
    \item Max 23 squad players per match per club.
    \item Red/yellow card suspension rules (e.g., 5 yellows = 1 match suspension).
    \item Numeric constraints such as \texttt{market\_value} $> 0$, \texttt{height\_cm} $> 0$, and \texttt{years\_experience} $\ge 0$ must hold.
    \item The foundation year of a club must be a valid past year.
    \item A player may not have two active contracts of the same type at the same time.
    \item A player may be associated with at most two clubs simultaneously (one permanent and one loan).
    \item The end date of a contract must be strictly later than its start date.
    \item Transfer constraints: Free transfer fee $= 0$; Purchase/Loan fee $> 0$; Loan transfer requires an active Permanent contract with parent club (and the parent club must match \texttt{from\_club\_id}).
    \item Manager/Club temporal exclusivity in Leads (a Manager cannot lead two Clubs simultaneously, and a Club cannot have two Managers simultaneously).
    \item A person's \texttt{date\_of\_birth} must be strictly in the past.
    \item A player cannot appear in the same match as both a home-team and away-team participant.
    \item 2 yellow cards in a single match automatically result in a red card (derived state).
    \item The Match Result (Win/Loss/Draw) is dynamically derived from home/away goals, rather than explicitly stored as an independent attribute.
\end{enumerate}

\end{document}
